\documentclass[12pt,a4paper]{article}
\usepackage{mathwizards-ingles}
\title{%
  \vspace{-1cm}
  {\Huge\bfseries\textcolor{mes2}{Mes 2 · Semana 8}}\\[0.3cm]
  {\Large Present Perfect, Experiencias y Transición Monolingüe}\\[0.2cm]
  {\large\color{grissuave} Nivel A2 · Hacia Definiciones en Inglés}\\[0.5cm]
  \rule{\textwidth}{1.5pt}
}
\author{}
\date{}

\begin{document}
\maketitle
\thispagestyle{empty}

\begin{cajanivel}
\textbf{Objetivo:} Se introduce el \textbf{Present Perfect} como puente entre pasado y presente. El estudiante aprende a hablar de \textbf{experiencias de vida} y \textbf{acciones recientes}. Además, esta semana marca la \textbf{transición hacia definiciones monolingües} (inglés-inglés) en Anki y Yomitan, reduciendo gradualmente la dependencia del español.
\end{cajanivel}

% ═══════════════════════════════════════════════════════════════════════════════
\section{Gramática Emergente: Present Perfect}
% ═══════════════════════════════════════════════════════════════════════════════

\begin{cajagrammar}[Present Perfect: \textit{have/has} + past participle]
Conecta el \textbf{pasado con el presente}. No importa \textit{cuándo} ocurrió, sino que tiene relevancia \textit{ahora}.

\medskip
\begin{tabularx}{\textwidth}{@{} l X @{}}
\toprule
\textbf{Estructura} & \textbf{Ejemplo} \\
\midrule
\textbf{Afirmativo:} Sujeto + have/has + V3 & \textit{I \textbf{have visited} Paris.} / \textit{She \textbf{has finished} her homework.} \\
\textbf{Negativo:} Sujeto + haven't/hasn't + V3 & \textit{I \textbf{haven't seen} that movie.} \\
\textbf{Pregunta:} Have/Has + sujeto + V3? & \textit{\textbf{Have} you \textbf{eaten} sushi?} \\
\bottomrule
\end{tabularx}

\medskip
\begin{cajaejemplo}[El participio pasado (V3 / past participle)]
\textbf{Verbos regulares:} el participio es igual al pasado simple (V + \textbf{-ed}):\\
\textit{walk $\rightarrow$ walked $\rightarrow$ \textbf{walked}} \quad \textit{study $\rightarrow$ studied $\rightarrow$ \textbf{studied}}

\medskip
\textbf{Verbos irregulares:} tienen una tercera columna propia:

\begin{multicols}{2}
\begin{tabularx}{\linewidth}{@{} l l l @{}}
Base & Past & \textbf{Participle} \\
\midrule
go & went & \eng{gone} \\
see & saw & \eng{seen} \\
eat & ate & \eng{eaten} \\
do & did & \eng{done} \\
have & had & \eng{had} \\
\end{tabularx}

\columnbreak

\begin{tabularx}{\linewidth}{@{} l l l @{}}
Base & Past & \textbf{Participle} \\
\midrule
be & was/were & \eng{been} \\
take & took & \eng{taken} \\
give & gave & \eng{given} \\
write & wrote & \eng{written} \\
know & knew & \eng{known} \\
\end{tabularx}
\end{multicols}
\end{cajaejemplo}
\end{cajagrammar}

% ═══════════════════════════════════════════════════════════════════════════════
\section{Chunks: Marcadores del Present Perfect}
% ═══════════════════════════════════════════════════════════════════════════════

\begin{cajachunk}[Signal Words — Palabras Clave del Present Perfect]
\begin{tabularx}{\textwidth}{@{} l l X @{}}
\toprule
\textbf{Palabra} & \textbf{Significado} & \textbf{Ejemplo} \\
\midrule
\eng{ever} & alguna vez (preguntas) & \textit{Have you \textbf{ever} been to London?} \\
\eng{never} & nunca (afirmativo) & \textit{I have \textbf{never} eaten sushi.} \\
\eng{already} & ya (antes de lo esperado) & \textit{She has \textbf{already} finished.} \\
\eng{yet} & todavía / ya (neg. y preg.) & \textit{I haven't done it \textbf{yet}.} / \textit{Have you finished \textbf{yet}?} \\
\eng{just} & acabar de & \textit{I have \textbf{just} arrived.} \\
\eng{for} & durante (duración) & \textit{I have lived here \textbf{for} five years.} \\
\eng{since} & desde (punto de inicio) & \textit{She has worked here \textbf{since} 2020.} \\
\bottomrule
\end{tabularx}

\medskip
\begin{cajaejemplo}[Posición de los marcadores]
\begin{itemize}[nosep, leftmargin=0.8cm]
  \item \textbf{ever / never / already / just} van \textbf{entre have/has y el participio}:\\
    \textit{I have \textbf{never} \textbf{seen} that.}
  \item \textbf{yet} va \textbf{al final} de la oración:\\
    \textit{Have you finished \textbf{yet}?}
  \item \textbf{for / since} van \textbf{al final} con la expresión de tiempo:\\
    \textit{I have lived here \textbf{for three years}.}
\end{itemize}
\end{cajaejemplo}
\end{cajachunk}

% ═══════════════════════════════════════════════════════════════════════════════
\section{Chunks: Hablar de Experiencias}
% ═══════════════════════════════════════════════════════════════════════════════

\begin{cajachunk}[Experience Chunks — Hablar de lo que Has Hecho (o No)]
\begin{tabularx}{\textwidth}{@{} X l @{}}
\eng{I have been to \underline{\hspace{2cm}}.} & \esp{He estado en \underline{\hspace{2cm}}.} \\
\eng{Have you ever \underline{\hspace{2cm}}?} & \esp{¿Alguna vez has \underline{\hspace{2cm}}?} \\
\eng{I have never tried \underline{\hspace{2cm}}.} & \esp{Nunca he probado \underline{\hspace{2cm}}.} \\
\eng{She has already \underline{\hspace{2cm}}.} & \esp{Ella ya ha \underline{\hspace{2cm}}.} \\
\eng{We haven't \underline{\hspace{2cm}} yet.} & \esp{Todavía no hemos \underline{\hspace{2cm}}.} \\
\eng{How long have you \underline{\hspace{2cm}}?} & \esp{¿Cuánto tiempo llevas \underline{\hspace{2cm}}?} \\
\eng{I've lived here since \underline{\hspace{2cm}}.} & \esp{Vivo aquí desde \underline{\hspace{2cm}}.} \\
\eng{I've known him for \underline{\hspace{2cm}}.} & \esp{Lo conozco desde hace \underline{\hspace{2cm}}.} \\
\end{tabularx}
\end{cajachunk}

\begin{cajagrammar}[Contraste: Present Perfect vs. Past Simple]
\begin{tabularx}{\textwidth}{@{} l X X @{}}
\toprule
& \textbf{Present Perfect} & \textbf{Past Simple} \\
\midrule
\textbf{Cuándo} & Tiempo no específico / relevante ahora & Tiempo específico / terminado \\
\textbf{Señales} & ever, never, already, yet, just & yesterday, last week, in 2020, ago \\
\textbf{Ejemplo} & \textit{I \textbf{have been} to Paris.} & \textit{I \textbf{went} to Paris \textbf{in 2019}.} \\
\textbf{Foco} & La experiencia importa & El momento importa \\
\bottomrule
\end{tabularx}

\medskip
\begin{cajaejemplo}[Regla Práctica]
Si puedes responder ``¿cuándo?'' con una fecha/momento concreto $\rightarrow$ \textbf{Past Simple}.\\
Si no importa el cuándo, solo importa si ocurrió $\rightarrow$ \textbf{Present Perfect}.
\end{cajaejemplo}
\end{cajagrammar}

% ═══════════════════════════════════════════════════════════════════════════════
\section{Gramática Emergente: Relative Clauses (\textit{who, which, that})}
% ═══════════════════════════════════════════════════════════════════════════════

\begin{cajagrammar}[Defining Relative Clauses — Cláusulas Relativas Definitorias]
Sirven para \textbf{dar información esencial} sobre un sustantivo, sin necesidad de una nueva oración.

\medskip
\begin{tabularx}{\textwidth}{@{} l l X @{}}
\toprule
\textbf{Pronombre} & \textbf{Se refiere a} & \textbf{Ejemplo} \\
\midrule
\eng{who} & personas & \textit{The woman \textbf{who} lives next door is a teacher.} \\
\eng{which} & cosas / animales & \textit{The book \textbf{which} I bought is very good.} \\
\eng{that} & personas o cosas & \textit{The movie \textbf{that} we watched was funny.} \\
\eng{where} & lugares & \textit{The restaurant \textbf{where} we ate was expensive.} \\
\bottomrule
\end{tabularx}

\medskip
\begin{cajaejemplo}[Nota]
En el inglés hablado cotidiano, \textbf{that} reemplaza a \textit{who} y \textit{which} en la mayoría de casos:\\
\textit{The man \textbf{that} called you...} = \textit{The man \textbf{who} called you...}
\end{cajaejemplo}
\end{cajagrammar}

% ═══════════════════════════════════════════════════════════════════════════════
\section{Transición Monolingüe}
% ═══════════════════════════════════════════════════════════════════════════════

\begin{cajatransicion}[Guía: De Definiciones Bilingües a Monolingües]
A partir de esta semana, se inicia la transición gradual para reducir la dependencia del español:

\medskip
\textbf{¿Por qué cambiar?} Las definiciones en inglés simple obligan al cerebro a \textbf{pensar en inglés}, no a traducir. Esto acelera dramáticamente la adquisición.

\medskip
\textbf{Pasos de la transición:}
\begin{enumerate}[leftmargin=1.2cm]
  \item \textbf{Semana 8:} Cambiar el diccionario principal de Yomitan a un \textit{Learner's Dictionary} monolingüe (Oxford/Cambridge para aprendices). Mantener el bilingüe como \textit{backup}.
  
  \item \textbf{Semanas 9-10:} En las tarjetas Anki nuevas, usar la definición en inglés en el reverso. Si no se entiende la definición, agregar la traducción \textbf{debajo} como apoyo secundario.
  
  \item \textbf{Semanas 11-12:} Eliminar la traducción al español por completo en las tarjetas nuevas. Las tarjetas antiguas con traducción se mantienen tal cual.
\end{enumerate}

\medskip
\textbf{Ejemplo de tarjeta monolingüe:}

\begin{center}
\begin{tabularx}{0.85\textwidth}{|X|}
\hline
\textbf{FRENTE:} \\
\textit{I have never \underline{\textbf{tried}} sushi, but I want to.} \\
\hline
\textbf{REVERSO:} \\
\textbf{tried} /traɪd/ — past tense of \textit{try}. \\
\textit{To taste or use something for the first time to see if you like it.} \\
\hline
\end{tabularx}
\end{center}
\end{cajatransicion}

% ═══════════════════════════════════════════════════════════════════════════════
\section{Oraciones Listas para Minar}
% ═══════════════════════════════════════════════════════════════════════════════

\begin{cajamineria}[Oraciones \textit{i+1} — Semana 8]

\begin{enumerate}[leftmargin=1.2cm]
  \item \textit{Have you \underline{ever} been to a different country?}\\
    {\small\textbf{Objetivo:} el marcador ``ever'' en preguntas de experiencia.}

  \item \textit{I have \underline{never} tried Indian food.}\\
    {\small\textbf{Objetivo:} ``never'' como negación dentro del Present Perfect.}

  \item \textit{She has \underline{already} finished her project.}\\
    {\small\textbf{Objetivo:} ``already'' para acciones completadas antes de lo esperado.}

  \item \textit{We haven't decided \underline{yet}.}\\
    {\small\textbf{Objetivo:} ``yet'' al final en oraciones negativas.}

  \item \textit{I've \underline{just} arrived. Give me a minute.}\\
    {\small\textbf{Objetivo:} ``just'' para acciones recientes.}

  \item \textit{How long have you \underline{lived} in this city?}\\
    {\small\textbf{Objetivo:} estructura ``How long have you + V3'' para duración.}

  \item \textit{I've known her \underline{since} we were children.}\\
    {\small\textbf{Objetivo:} ``since'' para punto de inicio en el tiempo.}

  \item \textit{The man \underline{who} called you is my brother.}\\
    {\small\textbf{Objetivo:} pronombre relativo ``who'' para personas.}

  \item \textit{This is the restaurant \underline{where} we had dinner last year.}\\
    {\small\textbf{Objetivo:} pronombre relativo ``where'' para lugares.}

  \item \textit{I have \underline{been} to three concerts this year.}\\
    {\small\textbf{Objetivo:} la expresión ``have been to'' (haber ido/estado en).}
\end{enumerate}
\end{cajamineria}

\vfill
\begin{center}
\rule{0.5\textwidth}{0.5pt}\\[0.3cm]
{\small\color{grissuave} Curso de Inmersión en Inglés · Mes 2 · Semana 8 · Nivel A2}
\end{center}

\end{document}
